\documentclass[12pt]{article}
\usepackage[utf8]{inputenc}
\usepackage[T1]{fontenc}
\usepackage{hyperref}
\title{Xingoni Language Resources for K-12 Curriculum Development}
\author{Compiled by Claude Sonnet 5.5 (claude-sonnet-5-5) for Andr\'as Kornai}
\date{}
\begin{document}
\maketitle

\section{Introduction}
Xingoni (ISO 639-3 xnq; Glottocode moza1251) is a Bantu language of northern Mozambique, in the Manda-Ngoni group. It is spoken mainly in Cabo Delgado and Niassa provinces. It is not a Nguni language and only shares the name. It is only marginally mutually intelligible with the Ngoni of Tanzania, and speakers often use Swahili to communicate with others. Documentation found is limited to a grammar sketch.

\section{Primary Educational Resources}
No primers, textbooks or school programmes were found in search results.

\section{Linguistic and Reference Materials}
\begin{itemize}
\item SIL Mozambique, \emph{Algumas notas gramaticais sobre Xingoni} (2008), a grammar sketch in Portuguese. I did not read it: \url{https://publikationen.ub.uni-frankfurt.de/opus4/files/22301/SILMozambique_GramaticaisXingoni.pdf}
\item Ngoni language (Mozambique) (Wikipedia): \url{https://en.wikipedia.org/wiki/Ngoni_language_(Mozambique)}
\end{itemize}

\section{Religious and Community Resources}
\begin{itemize}
\item Joshua Project profile of the Xingoni in Mozambique (demographic and missions information): \url{https://joshuaproject.net/index.php/people_groups/21982}
\end{itemize}

\section{Digital and Audio Resources}
\begin{itemize}
\item Omniglot page on Ngoni (Chingoni/Xingoni). Its sample text and numbers are for the Tanzanian variety, not Mozambican Xingoni: \url{https://omniglot.com/writing/ngoni.htm}
\end{itemize}

\section{Curriculum Development Context}
The language is spoken as a first language but is not known to be taught in schools. The only substantial material found is the SIL grammar sketch, so a curriculum would first need a vocabulary, an orthography and reading materials. These would be developed with speaker communities and Mozambican education authorities. Tanzanian Ngoni materials should not be assumed to apply.
\end{document}